% Source: https://artofproblemsolving.com/community/c3h1570419_distance_problem
% Stable identity: aops-topic-1570419-q01
Since $1+2+\cdots+19=\dfrac{19\cdot20}{2}=190$, she has walked exactly $19$ legs, of lengths $1,2,\dots,19$, in the directions north, east, south, west repeating. Put home at the origin with north the positive $y$-direction.

\textbf{North-south.} The odd legs alternate north and south:
\[
y=1-3+5-7+9-11+13-15+17-19=-10 .
\]
\textbf{East-west.} The even legs alternate east and west:
\[
x=2-4+6-8+10-12+14-16+18=10 .
\]

\begin{asy}
usepackage("amsmath");
size(5.5cm);
defaultpen(fontsize(10pt));
pair p=(0,0); pair[] d={(0,1),(1,0),(0,-1),(-1,0)};
path g=p;
for(int k=1;k<=19;++k){p=p+k*d[(k-1)%4]; g=g--p;}
draw(g,blue+linewidth(0.6));
draw((0,0)--p,red+linewidth(1.1));
dot((0,0)); dot(p);
label("$O$",(0,0),W);label("$(10,-10)$",p,SE);
label("$10\sqrt{2}$",p+(0.3,0.9),E,red);
label("\mbox{}",(min(currentpicture,user=true)+max(currentpicture,user=true))/2,black); // keeps the text after the figure black when LaTeX typesets the labels inline
\end{asy}

(Home is $O$; the blue path is her walk.) So Brenda ends at $(10,-10)$, and her distance from home is
\[
\sqrt{10^2+10^2}=\boxed{10\sqrt2}\text{ blocks}.
\]
